% =====================================================================
%  EXTENDED ABSTRACT (internal review only)
%  IEEE Access limits the abstract to 250 words, and its title block cannot
%  break across pages, so the boxed abstract in main.tex stays the
%  submission abstract. This section expands it for readers of the draft
%  and is switched off for submission with \extendedabstractfalse.
%  Every number repeats Sections V and VI and must stay identical to them.
% =====================================================================
\section*{Extended Abstract}
\label{sec:extended-abstract}

\emph{Problem.} Services that sell language models in tiers ask the user
to choose a model size and, increasingly, whether the model should reason
before it answers. Users choose badly in two opposite directions: some start
small and move up after every wrong answer, paying in full for each
discarded attempt and for resending the prompt, while others leave a large
reasoning configuration selected from an earlier task and pay its premium on
prompts that never needed it. Routing studies compare routers with fixed
single-model policies, and none that we located measures what this trial and
error costs.

\emph{Construct.} We define the \emph{escalation tax} as the cost a policy
pays beyond the cheapest configuration that would have answered the prompt,
the minimum sufficient configuration $\cstar(q)$. The tax splits into
\emph{escalation}, the cost of cheaper attempts that failed, and
\emph{stickiness}, the excess of staying on a configuration costlier than
$\cstar(q)$. We ask whether a small router placed in front of the
configurations can select $\cstar(q)$ in one shot, before any generation,
and whether its saving survives once the router pays for its own tokens,
latency and energy.

\emph{Setup.} The local grid crosses four Qwen3 sizes (0.6B, 1.7B, 4B and
8B, quantised to Q4\_K\_M) with thinking on and off, eight configurations
served on one laptop RTX~4060 GPU whose board power is sampled every
100\,ms. Every prompt runs in every configuration: 1{,}200 gradable prompts
from mathematics, code and multiple choice, and 400 open-ended requests
graded by a model judge. The gradable grid alone is 9{,}600 generations and
31.6 hours of GPU time. We price four manual escalation policies, train
classical, encoder and language-model routers under net-of-router
accounting, and compare them with a self-verifying cascade in the style of
AutoMix. Two further pools extend the result beyond the laptop: thirteen
public models priced by their providers, and six hosted configurations that
cross three GPT-5 sizes with two reasoning efforts.

\emph{The tax is large, and most of it is stickiness.} Leaving the largest
thinking configuration selected pays 12.96 times the minimum sufficient
energy and solves 79.8\% of prompts, because that configuration is not the
most accurate: between 16\% and 20\% of thinking generations reach the
4{,}096-token cap and fail, and 8B without thinking answers 3.8 points more
for an eighth of the energy. The frugal climb, which moves up sizes with
thinking off and turns thinking on only at the top, pays 0.89 times and
solves 92.2\%, at 945.7\,J per solved prompt against the oracle's 634.8\,J,
because its first four rungs are the four cheapest configurations and
together they are $\cstar(q)$ for 89.2\% of prompts. Turning thinking on
before moving up a size pays 8.66 times. The same ordering holds on the
public pool, at 7.96 against 2.69 times, and on the hosted ladder, at 4.35
against 0.22.

\emph{Reasoning has its own energy signature.} Thinking costs 6.4 to 8.3
times the energy per solved prompt of the same size without it, and the
ratio depends more on the task family than on size. Below 8B it was never
the cheaper way to buy accuracy: the next size up without thinking answered
more prompts for 3.8 to 5.5 times less energy per solved prompt. Rerunning a
sample of capped generations with an 8{,}192-token budget rescues 32.4\% of
them.

\emph{The unit changes the verdict.} Within a size, output tokens price the
mode axis to within about 15\%, but an output token costs 6.6 times as much
energy at 8B as at 0.6B, which a token count cannot see. Counted in tokens,
sticking with the largest model looks cheaper than trying thinking first,
and routing one tier too low looks cheaper than routing one too high. The
meter reverses both: in energy, routing too low costs 2.71 times as much as
routing too high.

\emph{Whether a router beats a disciplined user depends on the recovery
rule.} On 243 held-out prompts, with a failed route resent to the largest
thinking configuration, the best router, logistic regression on TF-IDF
features, spends 1.30 times the frugal climb's energy per solved prompt and
answers 3.7 points fewer. That configuration is expensive and not the most
accurate; when failures climb from the routed configuration instead, every
router comes within 6\% of the climb and four are 2.6\% to 3.4\% cheaper
than it at the same solved rate. Detection matters as much: when users accept
15\% of wrong answers as correct, the router answers 7.1 points more than the
climb, because it makes fewer attempts at which to be fooled. Against habits
the picture reverses outright: per prompt, the same router spends 61\% less
than leaving the largest configuration selected, 29\% less than starting at a
random configuration, and 34\% less than the cascade, whose self-check
accepts a wrong answer on a third of answerable prompts. Router overhead
stays under 1\% of the energy it routes, so
charging for it barely changes the ranking in energy; in tokens, the
language-model routers pay for their own prompts and fall behind.

\emph{Open-ended requests differ, with a caveat.} The judge agrees with a
human at $\kappa = 0.42$, inside the caveat band fixed before labelling, and
it is lenient. On these requests accuracy rises with both size and
thinking, from 46.3\% to 86.3\%, the cap barely binds, thinking costs only
about twice the energy per solved prompt, and at 0.6B it is cheaper than
moving up a size. The frugal climb still pays the least tax, 1.46 times,
against 5.39 for the sticky default.

\emph{Beyond the laptop.} On the public pool only four of thirteen models lie
on the cost--accuracy frontier, and the oracle pays \$25.47 per thousand
solved queries against \$102.08 for sticking with the top model. On the
hosted ladder high reasoning effort costs 5.45 to 11.8 times minimal effort,
and \texttt{gpt-5} at high effort answers fewer queries than
\texttt{gpt-5-mini} at high effort at 5.6 times its price. Pairing the hosted
and local labels prompt by prompt shows that the choice of reasoning mode
transfers between pools while the choice of size does not.

\emph{Implications and limits.} Defaults decide most of the tax: a service
that opens on a small configuration with reasoning off puts its users on the
frugal climb without any router, and for a user the evidence gives a short
rule: start small with reasoning off, move up a size before turning reasoning
on, and choose reasoning per prompt. A one-shot router earns its place against
costly defaults and for users who cannot check an answer; against a careful
user it needs a better recovery rule. The manual policies are modelled rather
than observed, the pool is one quantised family up to 8B on one machine, the
energy boundary is the GPU board, and the open-ended labels rest on a judge
of moderate agreement.
